\chapter{Raciocínio Lógico}

\section{Fundamentos da lógica proposicional}

A lógica proposicional estuda afirmações às quais se pode atribuir um valor lógico. O ponto de partida é a \termo{proposição}: uma sentença declarativa que pode ser classificada como verdadeira ou falsa. Não é necessário saber, de imediato, qual dos dois valores ela possui; basta que faça sentido atribuir-lhe um deles.

Assim, ``O TCE é órgão de controle externo'' é uma proposição. Já ``Feche a porta'', ``Qual é o resultado?'' e ``$x+2=7$'' não são proposições completas: a primeira é uma ordem, a segunda é uma pergunta e a terceira depende do valor da variável $x$.

\begin{teoria}[proposição e sentença aberta]
Uma \textbf{sentença aberta} contém variável livre e somente se torna proposição depois que a variável é determinada ou quantificada. Por exemplo,
\[
P(x): x>10
\]
não possui valor lógico enquanto $x$ estiver livre. Já ``para todo $x\in\mathbb{R}$, se $x>10$, então $x>5$'' é uma proposição quantificada.
\end{teoria}

\subsection{Proposições simples e compostas}

Uma proposição é \termo{simples} quando não possui outra proposição como componente lógico. É \termo{composta} quando resulta da combinação de proposições simples por conectivos.

Considere:
\begin{itemize}
\item $p$: ``o relatório foi aprovado'';
\item $q$: ``o processo foi arquivado''.
\end{itemize}

A expressão ``o relatório foi aprovado e o processo foi arquivado'' corresponde a $p\land q$ e é uma proposição composta.

\section{Conectivos lógicos e tabelas-verdade}

Os conectivos fundamentais determinam o valor lógico da proposição composta a partir dos valores de suas partes.

\begin{table}[ht]
\centering
\small
\begin{tabularx}{\textwidth}{l c X}
\toprule
Operação & Símbolo & Regra de verdade \\
\midrule
Negação & $\neg p$ & inverte o valor lógico de $p$. \\
Conjunção & $p\land q$ & verdadeira somente quando $p$ e $q$ são verdadeiras. \\
Disjunção inclusiva & $p\lor q$ & falsa somente quando $p$ e $q$ são falsas. \\
Disjunção exclusiva & $p\oplus q$ & verdadeira quando exatamente uma das proposições é verdadeira. \\
Condicional & $p\to q$ & falsa somente quando $p$ é verdadeira e $q$ é falsa. \\
Bicondicional & $p\leftrightarrow q$ & verdadeira quando $p$ e $q$ possuem o mesmo valor lógico. \\
\bottomrule
\end{tabularx}
\end{table}

\subsection{Tabela-verdade completa}

\begin{center}
\begin{tabular}{cc|ccccc}
\toprule
$p$ & $q$ & $p\land q$ & $p\lor q$ & $p\oplus q$ & $p\to q$ & $p\leftrightarrow q$ \\
\midrule
V & V & V & V & F & V & V \\
V & F & F & V & V & F & F \\
F & V & F & V & V & V & F \\
F & F & F & F & F & V & V \\
\bottomrule
\end{tabular}
\end{center}

A disjunção representada por $\lor$ é, em regra, \textbf{inclusiva}: admite que as duas proposições sejam verdadeiras simultaneamente. Expressões como ``ou...ou'', quando o contexto exigir exclusividade, podem ser modeladas por $\oplus$.

\begin{exemplo}[disjunção inclusiva e exclusiva]
``O servidor possui graduação em Direito ou Administração'' pode ser verdadeira se ele possuir as duas graduações: leitura inclusiva.

``Ou o equipamento está ligado à rede A ou está ligado à rede B, mas não às duas'' explicita exclusividade: exatamente uma condição deve ocorrer.
\end{exemplo}

\subsection{Número de linhas da tabela-verdade}

Se uma proposição composta contém $n$ proposições simples logicamente independentes, a tabela-verdade possui
\[
2^n
\]
linhas. Com três proposições simples, são $2^3=8$ combinações; com cinco, $2^5=32$.

\subsection{Precedência dos conectivos}

Quando não há parênteses suficientes para eliminar ambiguidades, usa-se uma precedência convencional. Uma ordem didática comum é:
\[
\neg \quad > \quad \land \quad > \quad \lor \quad > \quad \to \quad > \quad \leftrightarrow.
\]

Parênteses sempre prevalecem. Em prova, a leitura correta da estrutura é mais importante que memorizar uma convenção isoladamente.

\section{A condicional em profundidade}

A proposição $p\to q$ é chamada condicional. $p$ é o \termo{antecedente} e $q$ é o \termo{consequente}. Ela somente é falsa quando o antecedente ocorre e o consequente não ocorre.

\subsection{Condição suficiente e condição necessária}

Em
\[
p\to q,
\]
$p$ é condição \textbf{suficiente} para $q$, enquanto $q$ é condição \textbf{necessária} para $p$.

Considere:

\begin{quote}
Se um processo foi julgado definitivamente, então existe decisão final registrada.
\end{quote}

``Ser julgado definitivamente'' é suficiente para concluir que existe decisão final registrada. A existência de decisão final é necessária para que a primeira condição seja verdadeira.

\begin{atencao}[somente se]
A expressão ``$p$ somente se $q$'' significa $p\to q$. Portanto, em ``o acesso é liberado somente se houver autenticação'', autenticação é condição necessária para liberação.
\end{atencao}

\subsection{Recíproca, inversa e contrapositiva}

A partir de $p\to q$:
\begin{itemize}
\item \textbf{recíproca}: $q\to p$;
\item \textbf{inversa}: $\neg p\to\neg q$;
\item \textbf{contrapositiva}: $\neg q\to\neg p$.
\end{itemize}

A proposição original é logicamente equivalente à contrapositiva:
\[
p\to q\equiv\neg q\to\neg p.
\]
A recíproca e a inversa são equivalentes entre si, mas não, em geral, à proposição original.

\begin{exemplo}[contrapositiva]
Se todo backup válido possui cópia íntegra, então, ao constatar que uma cópia não é íntegra, pode-se concluir pela contrapositiva que ela não constitui backup válido segundo a premissa adotada.

Não se pode concluir, porém, que toda cópia íntegra seja necessariamente um backup válido. Isso seria afirmar a recíproca.
\end{exemplo}

\section{Equivalências lógicas}

Duas proposições são logicamente equivalentes quando possuem o mesmo valor lógico em todas as valorações possíveis. As equivalências permitem simplificar expressões e resolver questões sem construir tabelas completas.

\subsection{Equivalências fundamentais}

\begin{align*}
p\to q &\equiv \neg p\lor q,\\
\neg(p\to q) &\equiv p\land\neg q,\\
p\leftrightarrow q &\equiv (p\to q)\land(q\to p),\\
p\leftrightarrow q &\equiv (p\land q)\lor(\neg p\land\neg q).
\end{align*}

\subsection{Leis de De Morgan}

\begin{align*}
\neg(p\land q) &\equiv \neg p\lor\neg q,\\
\neg(p\lor q) &\equiv \neg p\land\neg q.
\end{align*}

A negação ``entra'' na expressão trocando conjunção por disjunção, ou disjunção por conjunção, e negando cada componente.

\begin{exemplo}[negação correta]
A negação de ``o sistema está disponível e o banco está íntegro'' é:

``o sistema não está disponível \textbf{ou} o banco não está íntegro''.

Não é necessário que as duas condições falhem simultaneamente.
\end{exemplo}

\subsection{Outras leis úteis}

\begin{align*}
\neg\neg p &\equiv p,\\
p\land p &\equiv p,\\
p\lor p &\equiv p,\\
p\land q &\equiv q\land p,\\
p\lor q &\equiv q\lor p,\\
p\land(q\lor r)&\equiv(p\land q)\lor(p\land r),\\
p\lor(q\land r)&\equiv(p\lor q)\land(p\lor r).
\end{align*}

Também são importantes as leis de absorção:
\[
p\lor(p\land q)\equiv p,
\qquad
p\land(p\lor q)\equiv p.
\]

\section{Negação de estruturas frequentes}

A negação deve atuar sobre o \textbf{escopo completo} da afirmação.

\begin{table}[ht]
\centering
\small
\begin{tabularx}{\textwidth}{X X}
\toprule
Proposição & Negação \\
\midrule
$p\land q$ & $\neg p\lor\neg q$ \\
$p\lor q$ & $\neg p\land\neg q$ \\
$p\to q$ & $p\land\neg q$ \\
$p\leftrightarrow q$ & $(p\land\neg q)\lor(\neg p\land q)$ \\
Todo A é B & Existe pelo menos um A que não é B \\
Nenhum A é B & Existe pelo menos um A que é B \\
Algum A é B & Nenhum A é B \\
\bottomrule
\end{tabularx}
\end{table}

\begin{cebraspe}
``Nem todos os contratos foram auditados'' equivale a ``pelo menos um contrato não foi auditado''. Não equivale a ``nenhum contrato foi auditado''.
\end{cebraspe}

\section{Tautologia, contradição e contingência}

Uma \termo{tautologia} é verdadeira em todas as valorações. Uma \termo{contradição} é falsa em todas. Uma \termo{contingência} é verdadeira em algumas e falsa em outras.

Exemplos clássicos:
\[
p\lor\neg p \quad \text{é tautologia},
\]
\[
p\land\neg p \quad \text{é contradição}.
\]

A expressão
\[
(p\to q)\leftrightarrow(\neg q\to\neg p)
\]
é uma tautologia porque expressa a equivalência entre a condicional e sua contrapositiva.

\subsection{Técnica da busca de contraexemplo}

Para provar que uma expressão \textbf{não} é tautologia, basta encontrar uma valoração em que ela seja falsa. Para testar um argumento, basta procurar uma situação em que todas as premissas sejam verdadeiras e a conclusão falsa.

Essa técnica é especialmente eficiente em itens Certo/Errado: um único contraexemplo destrói uma afirmação universal.

\section{Argumentação lógica e validade}

Um argumento possui premissas e conclusão. Ele é \termo{válido} quando é impossível que todas as premissas sejam verdadeiras e a conclusão seja falsa simultaneamente.

Validade é propriedade da \textbf{forma lógica}; não significa que as premissas sejam factualmente verdadeiras.

\subsection{Regras clássicas de inferência}

\paragraph{Modus ponens}
\[
p\to q,\quad p \quad\therefore\quad q.
\]

\paragraph{Modus tollens}
\[
p\to q,\quad \neg q \quad\therefore\quad \neg p.
\]

\paragraph{Silogismo hipotético}
\[
p\to q,\quad q\to r \quad\therefore\quad p\to r.
\]

\paragraph{Silogismo disjuntivo}
\[
p\lor q,\quad \neg p \quad\therefore\quad q.
\]

\subsection{Falácias formais recorrentes}

\textbf{Afirmação do consequente}:
\[
p\to q,\quad q \quad\therefore\quad p.
\]
É inválida.

\textbf{Negação do antecedente}:
\[
p\to q,\quad \neg p \quad\therefore\quad \neg q.
\]
Também é inválida.

\begin{exemplo}[validade e verdade]
Premissas: ``Se todo auditor é médico, então Ana é médica'' e ``todo auditor é médico''. Mesmo que a premissa geral seja falsa no mundo real, pode-se estudar a validade formal do argumento independentemente de sua verdade factual.
\end{exemplo}

\section{Quantificadores e lógica de predicados}

A lógica de predicados permite formalizar afirmações sobre elementos de um universo.

\subsection{Quantificador universal}

``Todo A é B'' pode ser escrito como
\[
\forall x\,(A(x)\to B(x)).
\]
Observe o uso da condicional: para cada elemento, se ele pertence a A, então pertence a B.

\subsection{Quantificador existencial}

``Algum A é B'' corresponde a
\[
\exists x\,(A(x)\land B(x)).
\]
É necessário que exista pelo menos um elemento que satisfaça as duas propriedades.

\subsection{Negação de quantificadores}

\begin{align*}
\neg\forall x\,P(x)&\equiv\exists x\,\neg P(x),\\
\neg\exists x\,P(x)&\equiv\forall x\,\neg P(x).
\end{align*}

Em linguagem natural:
\begin{itemize}
\item negar ``todos possuem a característica'' produz ``pelo menos um não possui'';
\item negar ``existe alguém com a característica'' produz ``ninguém possui a característica''.
\end{itemize}

\section{Diagramas lógicos e relações entre conjuntos}

Proposições categóricas podem ser visualizadas por conjuntos.

\begin{itemize}
\item Todo A é B: $A\subseteq B$.
\item Nenhum A é B: $A\cap B=\varnothing$.
\item Algum A é B: $A\cap B\neq\varnothing$.
\item Algum A não é B: $A-B\neq\varnothing$.
\end{itemize}

\begin{exemplo}[conclusão necessária]
Premissas:
\begin{enumerate}
\item todo auditor de TI é servidor;
\item nenhum terceirizado é servidor.
\end{enumerate}

Como o conjunto dos auditores de TI está contido no conjunto dos servidores e o conjunto dos terceirizados é disjunto do conjunto dos servidores, conclui-se necessariamente que nenhum auditor de TI é terceirizado, dentro desse universo hipotético.
\end{exemplo}

Cuidado com existência. De ``todo A é B'' não se conclui, por si só, que exista algum A. A inclusão pode ser verdadeira mesmo se $A$ for vazio.

\section{Teoria dos conjuntos}

Um conjunto é uma coleção de elementos. As operações básicas são união, interseção, diferença e complemento.

\begin{align*}
A\cup B &= \{x: x\in A \text{ ou } x\in B\},\\
A\cap B &= \{x: x\in A \text{ e } x\in B\},\\
A-B &= \{x: x\in A \text{ e } x\notin B\}.
\end{align*}

\subsection{Princípio da inclusão-exclusão}

Para dois conjuntos finitos:
\[
|A\cup B|=|A|+|B|-|A\cap B|.
\]

Para três:
\[
|A\cup B\cup C|=|A|+|B|+|C|-|A\cap B|-|A\cap C|-|B\cap C|+|A\cap B\cap C|.
\]

\begin{exemplo}[dupla contagem]
Em um grupo de 120 servidores, 70 utilizam o sistema A, 55 utilizam o sistema B e 25 utilizam ambos. Logo,
\[
|A\cup B|=70+55-25=100.
\]
Portanto, 20 servidores não utilizam nenhum dos dois sistemas.
\end{exemplo}

\section{Razão, proporção e grandezas proporcionais}

A razão entre $a$ e $b$, com $b\neq0$, é $a/b$. Uma proporção é uma igualdade entre razões:
\[
\frac{a}{b}=\frac{c}{d}\quad\Longrightarrow\quad ad=bc.
\]

Duas grandezas são \textbf{diretamente proporcionais} quando multiplicar uma por um fator multiplica a outra pelo mesmo fator. São \textbf{inversamente proporcionais} quando o produto permanece constante.

\begin{exemplo}[proporção inversa]
Se 6 servidores realizam uma tarefa em 10 dias, mantendo produtividade constante, o total de trabalho é proporcional a $6\times10=60$ servidor-dias. Com 12 servidores, o tempo teórico será
\[
60/12=5\text{ dias}.
\]
\end{exemplo}

\section{Porcentagem e variações sucessivas}

Uma variação percentual deve ser tratada por fator multiplicativo. Aumento de $r\%$ corresponde a multiplicar por
\[
1+\frac{r}{100},
\]
e redução por
\[
1-\frac{r}{100}.
\]

\subsection{Variações sucessivas}

Percentuais sucessivos incidem sobre bases diferentes e, por isso, não devem ser somados mecanicamente.

\begin{exemplo}[aumento e redução]
Um valor de 100 aumenta 20\%:
\[
100\times1{,}20=120.
\]
Depois sofre redução de 20\%:
\[
120\times0{,}80=96.
\]
O resultado é 4\% menor que o valor inicial. Aumento de 20\% seguido de redução de 20\% não se anula.
\end{exemplo}

\subsection{Variação relativa}

Se uma grandeza passa de $V_0$ para $V_1$, a variação percentual é
\[
\frac{V_1-V_0}{V_0}\times100\%.
\]
O denominador é o valor de referência. Trocar a base altera o percentual.

\section{Equações e sistemas lineares}

Problemas verbais devem ser convertidos em relações algébricas.

Uma equação de primeiro grau pode ser escrita como
\[
ax+b=c,\qquad a\neq0.
\]

\begin{exemplo}[modelagem]
Um setor possui 32 servidores entre analistas e técnicos. Há 8 analistas a mais que técnicos. Se $a$ é o número de analistas e $t$ o de técnicos:
\[
\begin{cases}
a+t=32,\\
a=t+8.
\end{cases}
\]
Substituindo:
\[
(t+8)+t=32\Rightarrow2t=24\Rightarrow t=12,
\]
e, portanto, $a=20$.
\end{exemplo}

O resultado matemático precisa ser compatível com o contexto. Quantidades de pessoas, processos ou equipamentos geralmente exigem valores inteiros não negativos.

\section{Sequências e reconhecimento de padrões}

Em sequências, não basta encontrar uma regra que funcione para dois ou três termos; a regra precisa ser compatível com todo o conjunto apresentado. Progressões aritméticas possuem diferença constante; progressões geométricas possuem razão constante. Outros padrões podem alternar operações, combinar subsequências ou depender da posição.

\begin{exemplo}[progressão aritmética]
Na sequência $5,9,13,17,\ldots$, a diferença é 4. O termo geral é
\[
a_n=5+(n-1)4=4n+1.
\]
Logo, $a_{20}=81$.
\end{exemplo}

\begin{cebraspe}[como a banca torna a questão difícil]
Em nível mais alto, a dificuldade não está em calcular uma tabela-verdade elementar, mas em identificar corretamente o escopo da negação, a direção de uma condição necessária, a validade de uma inferência ou a base percentual. Uma única palavra — ``somente'', ``todos'', ``algum'', ``necessário'', ``suficiente'' — pode alterar completamente a estrutura lógica.
\end{cebraspe}

\section{Síntese conceitual}

\mapafuncional{Raciocínio Lógico}
{Linguagem natural $\rightarrow$ formalização $\rightarrow$ relações lógicas $\rightarrow$ conclusão válida}
{Proposições $\rightarrow$ conectivos $\rightarrow$ tabelas e equivalências.\\Argumentos $\rightarrow$ premissas $\rightarrow$ regras de inferência.\\Quantificadores $\rightarrow$ conjuntos e diagramas.\\Problemas quantitativos $\rightarrow$ razão, porcentagem, equações e sequências.}
{Condicional: antecedente suficiente, consequente necessário.\\Negação: preserve o escopo e aplique De Morgan.\\Validade: procure premissas verdadeiras com conclusão falsa.\\Conjuntos: elimine dupla contagem.\\Percentuais: use fatores multiplicativos.}
{Recíproca $\neq$ contrapositiva.\\Ou inclusivo $\neq$ ou exclusivo.\\Nem todos $\neq$ nenhum.\\Validade $\neq$ verdade factual.\\+20\% e -20\% $\neq$ zero.}
{Qual é o conectivo principal?\\Há quantificador?\\Qual é o escopo da negação?\\A conclusão é necessária ou apenas possível?\\Existe dupla contagem?\\Qual é a base do percentual?}

\begin{resumo}
O conteúdo de Raciocínio Lógico deve ser dominado em dois níveis: \textbf{formal}, pela manipulação de símbolos, tabelas e equivalências; e \textbf{semântico}, pela tradução precisa de expressões como ``somente se'', ``todo'', ``algum'', ``necessário'' e ``suficiente''. A maior parte das questões seletivas combina esses dois níveis.
\end{resumo}
